\section{一页判据}
\label{sec:onepage}

本文的结论散在六节里。这一节把它们收成一套可以直接照着走的顺序。
\textbf{如果只读一节,读这一节。}

\subsection{四步，不可颠倒}

\begin{figure*}[t]
\centering
\small
\setlength{\fboxsep}{4pt}
\begin{tabular}{@{}c@{}}
\fbox{\parbox{15.0cm}{\centering\textbf{1\quad 看负载形态}\quad
  命中率 / 输入输出比 / 轮间隔 / 上下文长度\\[1pt]
  \footnotesize 命中率 $>90\%$: 先做前缀缓存与路由粘性, 分离退居其次
  \quad$\vert$\quad 输入输出比 $<1$: 在窗口左侧, 别分离}}
\\[2pt] $\Downarrow$
\\[2pt]
\fbox{\parbox{15.0cm}{\centering\textbf{2\quad 定 SLO 形状}\quad
  哪个更紧（TTFT 还是 TPOT）/ 两者比值 / 达成率目标\\[1pt]
  \footnotesize TTFT 紧: 优先 prefill、增大 token 预算、P 侧多配
  \quad$\vert$\quad TPOT 紧: 优先 decode、减小预算、D 侧多配}}
\\[2pt] $\Downarrow$
\\[2pt]
\fbox{\parbox{15.0cm}{\centering\textbf{3\quad 查窗口}\quad
  两阶段最优配置是否显著不同 / 是否有高带宽 KV 通路 / 是否落在窗口内\\[1pt]
  \footnotesize 任一条不满足: 不要分离; 已分离的应回退}}
\\[2pt] $\Downarrow$
\\[2pt]
\fbox{\parbox{15.0cm}{\centering\textbf{4\quad 每类选一个}\quad
  数据面 / 执行面 / 控制面各选, 四类技巧里每类选一个\\[1pt]
  \footnotesize 前置条件不满足时开了是负收益: LPLB 无冗余副本只有开销、
  逐层传输在 KV 太小时被固定开销吃掉、DBO 低于阈值不拆}}
\end{tabular}
\caption{判断顺序。前三步任一步不通过, 后面的技巧都不必谈——
方向错了, 做得越细越糟。}
\label{fig:decision}
\end{figure*}

\noindent 顺序不能换, 因为每一步都能\textbf{让后面步骤的结论失效},
而不只是让它更精确: 第 1 步能让一整类技巧失去用武之地（命中率
$94.2\%$ 时搬得更快没有意义）; 第 2 步能翻转技巧的方向（同一 token
预算在两个 SLO 下反向）; 第 3 步能直接否决分离本身。
反过来先问「用什么技巧」, 只会得到一堆互相冲突的建议。

\subsection{六层的关系}

\begin{table*}[t]
\centering
\caption{六次切法各自回答的问题。它们不是六个并列章节, 而是一条链:
前四层是\emph{否决层}, 后两层是\emph{执行层}。}
\label{tab:layers}
\small
\setlength{\tabcolsep}{7pt}
\begin{tabular}{@{}p{0.75cm}p{2.99cm}p{5.08cm}p{4.33cm}p{1.79cm}@{}}
\toprule
层 & 切法 & 回答什么 & 能否决什么 & 节 \\
\midrule
1 & 负载形态 & 输入长什么样 & 整类技巧的用武之地 & \ref{sec:load} \\
2 & SLO 形状 & 要什么 & 技巧的方向 & \ref{sec:slo} \\
3 & 收益边界 & 值不值得做 & 分离本身 & \ref{sec:boundary} \\
4 & 中间结论 & —— & 三者都成立的区间 & 本节 \\
5 & 实现平面 & 栈里有哪几层 & 保证不漏 & \ref{sec:plane} \\
6 & 四类技巧 & 传输账怎么省 & 给出同类替代 & \ref{sec:technique} \\
7 & 模型差异 & 具体模型怎么做 & 给出起点配置 & \ref{sec:model} \\
\bottomrule
\end{tabular}
\end{table*}

\noindent 常见错误是把六层当清单逐条做——那会做很多本不该做的事。
正确用法是\textbf{先跑否决层, 任一层否决就停}, 通过了再进执行层。

\subsection{本文没有回答什么}

一页判据看起来太像完整答案, 故这几条必须并列:

\begin{itemize}[leftmargin=1.3em,itemsep=1pt,topsep=2pt]
  \item \textbf{没有端到端复现。} 未在任何硬件上跑过分离部署,
        硬件相关数字都是文献值。
  \item \textbf{没有配比公式。} 只说配比应由两阶段的实测吞吐与延迟
        曲线决定, 没有闭式解。
  \item \textbf{窗口形状来自模拟。} 8 个模型的扫描全部来自同一套 NPU
        模拟设计空间, 窗口的\emph{位置}可能迁移, 但「跨一个数量级」
        这个\emph{形状}更可能稳。
  \item \textbf{一页判据不能替代测量。} 它只能排除错误方向,
        不能告诉你某个具体部署能快多少。
\end{itemize}

\noindent 若只能记一句: \textbf{分离是「让两阶段不必做同一件事」这条
主线里最贵的一个实例}。考虑它之前, 先看有没有更便宜的实例
（前缀命中、chunked prefill、按阶段配置）——往往有。
